% ============================================================
%  RAPPORT TECHNIQUE — TalentMatch
%  Youssef Gara — Stage PFE 3S — 2026
% ============================================================
\documentclass[12pt,a4paper]{report}

\usepackage[utf8]{inputenc}
\usepackage[T1]{fontenc}
\usepackage[french]{babel}
\usepackage{geometry}
\usepackage{setspace}
\usepackage{titlesec}
\usepackage{float}
\usepackage{booktabs}
\usepackage{array}
\usepackage{tabularx}
\usepackage{graphicx}
\usepackage{xcolor}
\usepackage{fancyhdr}
\usepackage{listings}
\usepackage{amsmath}
\usepackage{framed}

\geometry{top=2.5cm, bottom=2.5cm, left=3cm, right=2.5cm}
\onehalfspacing

% ── Couleurs ──────────────────────────────────────────────
\definecolor{navyblue}{RGB}{10,35,80}
\definecolor{goldcolor}{RGB}{184,142,76}
\definecolor{lightgray}{RGB}{245,245,245}
\definecolor{codegreen}{rgb}{0.1,0.5,0.1}
\definecolor{codegray}{rgb}{0.5,0.5,0.5}
\definecolor{codeblue}{rgb}{0.1,0.1,0.8}

% ── En-têtes / pieds de page ───────────────────────────────
\pagestyle{fancy}
\fancyhf{}
\fancyhead[L]{\small\textcolor{navyblue}{\textbf{TalentMatch — Rapport Technique}}}
\fancyhead[R]{\small\textcolor{goldcolor}{Youssef Gara — 3S — 2026}}
\fancyfoot[C]{\thepage}
\renewcommand{\headrulewidth}{0.4pt}

% ── Titre de chapitres ─────────────────────────────────────
\titleformat{\chapter}[display]
  {\normalfont\Large\bfseries\color{navyblue}}
  {\chaptertitlename\ \thechapter}{10pt}{\huge}
\titlespacing*{\chapter}{0pt}{-10pt}{20pt}

\titleformat{\section}{\bfseries\large\color{navyblue}}{\thesection}{1em}{}
\titleformat{\subsection}{\bfseries\normalsize\color{navyblue}}{\thesubsection}{1em}{}

% ── Box colorées (framed) ─────────────────────────────────
\newenvironment{infobox}[1][]{%
  \begin{framed}\noindent\textbf{\textcolor{navyblue}{#1}}\\[2pt]%
}{\end{framed}}

\newenvironment{codebox}[1][]{%
  \begin{framed}\noindent\textbf{\small\textcolor{codeblue}{#1}}\\[2pt]\ttfamily\small%
}{\end{framed}}

\newenvironment{whybox}[1][Pourquoi ce choix ?]{%
  \colorbox{yellow!20}{\textbf{\textcolor{navyblue}{#1}}}\vspace{2pt}
  \begin{framed}%
}{\end{framed}}

% ── Listings (code) ────────────────────────────────────────
\lstset{
  backgroundcolor=\color{lightgray},
  basicstyle=\ttfamily\small,
  keywordstyle=\color{codeblue}\bfseries,
  commentstyle=\color{codegreen}\itshape,
  stringstyle=\color{red!70},
  breaklines=true,
  frame=single,
  rulecolor=\color{gray!50},
  numbers=left, numberstyle=\tiny\color{gray},
  tabsize=4
}

% ── Tableau path ──────────────────────────────────────────
\newcommand{\filepath}[1]{\texttt{\small\textcolor{codeblue}{#1}}}

% ============================================================
\begin{document}
% ============================================================

% ── PAGE DE GARDE ─────────────────────────────────────────
\begin{titlepage}
\thispagestyle{empty}
\begin{center}
  \vspace*{1cm}
  {\color{navyblue}\rule{\textwidth}{2pt}}\\[0.5cm]
  {\Huge\bfseries\color{navyblue} TalentMatch}\\[0.3cm]
  {\Large Plateforme intelligente de matching CV / Offres d'emploi}\\[0.3cm]
  {\color{goldcolor}\rule{\textwidth}{1pt}}\\[1.5cm]

  {\LARGE\bfseries Rapport Technique}\\[0.4cm]
  {\large Choix technologiques, architecture et implémentation}\\[2cm]

  \begin{tabular}{ll}
    \textbf{Réalisé par :}     & Youssef Gara \\[4pt]
    \textbf{Entreprise :}      & 3S Group — Tunis \\[4pt]
    \textbf{Encadrante société :}  & Mme Sirine Nasri \\[4pt]
    \textbf{Encadrante académique :} & Mme Loujayne Bouzarti (Esprit) \\[4pt]
    \textbf{Période :}         & Février 2026 — Août 2026 \\[4pt]
    \textbf{Date :}            & Juin 2026 \\
  \end{tabular}\\[3cm]

  {\color{navyblue}\rule{\textwidth}{2pt}}
\end{center}
\end{titlepage}

\tableofcontents
\newpage

% ============================================================
\chapter{Contexte et objectif du projet}
% ============================================================

\section{Problématique}

Le recrutement moderne souffre d'un paradoxe : les entreprises reçoivent un
volume croissant de candidatures, mais les outils existants (LinkedIn, Indeed,
ANAPEC) se limitent à une recherche par \textbf{mots-clés}, sans comprendre
le sens réel des profils.

\begin{infobox}[Chiffre clé]
Un recruteur consulte un CV en moyenne \textbf{23 secondes}. Sans outil
intelligent, de nombreux bons profils sont ignorés ou mal classés.
\end{infobox}

Les problèmes concrets identifiés :
\begin{itemize}
  \item Faux positifs : un profil qui cite un mot-clé sans maîtriser le domaine.
  \item Bons profils ignorés : un expert qui formule ses compétences différemment.
  \item Pas de compréhension du \textbf{contexte sémantique} ni de la trajectoire professionnelle.
  \item Absence d'explication : le recruteur ne sait pas \textit{pourquoi} un candidat est classé 1\textsuperscript{er}.
\end{itemize}

\section{Objectif — TalentMatch}

TalentMatch est une plateforme web complète qui automatise le \textbf{matching CV / offres d'emploi}
via un pipeline IA :

\begin{enumerate}
  \item \textbf{Extraction} — lire un CV (PDF, DOCX, scan) et en extraire le texte brut.
  \item \textbf{Parsing NLP} — structurer le texte en JSON (compétences, expériences, formations).
  \item \textbf{Scoring IA} — calculer un score de pertinence 0–100\% par modèle d'embeddings.
  \item \textbf{Classement et explication} — classer les candidats et afficher un rapport détaillé.
\end{enumerate}

% ============================================================
\chapter{Architecture générale du système}
% ============================================================

\section{Vue d'ensemble}

Le système suit une architecture \textbf{client-serveur découplée} :

\begin{center}
\begin{tabular}{|c|c|c|}
\hline
\textbf{Couche} & \textbf{Technologie} & \textbf{Rôle} \\
\hline
Frontend & React 18 + Vite & Interface utilisateur (SPA) \\
API & FastAPI (Python) & Logique métier, endpoints REST \\
IA / NLP & spaCy, BGE-M3, MLP & Parsing et matching \\
Base de données & PostgreSQL 16 & Stockage persistant \\
Auth & JWT (python-jose) & Sécurité des accès \\
Email & SMTP Gmail & Notifications \\
Déploiement & Railway + Docker & Cloud \\
\hline
\end{tabular}
\end{center}

\section{Structure du code}

\begin{codebox}[Arborescence principale — projet 3s-talentmatch/]
backend/
  app/
    main.py           <- Point d'entrée FastAPI (lifespan, CORS, routes)
    routes/           <- Endpoints REST (auth, cv, matching, job_offers...)
    models/           <- Tables SQLAlchemy (ORM)
    schemas/          <- Validations Pydantic
    services/
      extraction/     <- Lecture PDF/DOCX/OCR
      nlp/            <- Pipeline NLP (spaCy)
      matching/       <- Moteur de matching principal
      matching_sandbox/ <- Scoreurs IA (BGE-M3, MLP)
      ai_summary/     <- Résumé IA du profil (Claude)
  scripts/            <- Scripts d'entraînement et de seed

frontend/
  src/
    pages/            <- Pages React (Dashboard, OfferDetail...)
    components/       <- Composants réutilisables (MatchReport...)
    services/api.js   <- Instance Axios centralisée
\end{codebox}

% ============================================================
\chapter{Choix technologiques et justifications}
% ============================================================

\section{Backend — FastAPI}

\begin{whybox}[Pourquoi FastAPI et pas Flask ou Django ?]
\begin{itemize}
  \item \textbf{Async natif} : FastAPI est basé sur ASGI (Starlette), ce qui permet
        de gérer plusieurs requêtes simultanées sans bloquer — crucial pour les
        appels aux modèles IA qui sont lents.
  \item \textbf{Documentation automatique} : Swagger UI sur \texttt{/docs} et
        ReDoc sur \texttt{/redoc} générés automatiquement depuis les types Python.
  \item \textbf{Validation automatique} : les schémas Pydantic valident les entrées/sorties.
  \item \textbf{Flask} : trop bas niveau, pas de validation intégrée, sync par défaut.
  \item \textbf{Django} : monolithique, overhead inutile pour une API pure.
\end{itemize}
\end{whybox}

\begin{infobox}[Fichier clé]
  \filepath{backend/app/main.py} — point d'entrée FastAPI.
  Le \textbf{lifespan handler} (lignes 1–50) pré-chauffe le modèle BGE-M3 au démarrage
  pour éviter la latence au premier appel de matching.
\end{infobox}

\section{Base de données — PostgreSQL}

\begin{whybox}[Pourquoi PostgreSQL et pas MySQL ou MongoDB ?]
\begin{itemize}
  \item \textbf{ACID} : transactions fiables — important pour les opérations de matching
        qui mettent à jour plusieurs tables en même temps.
  \item \textbf{JSON natif} : le champ \texttt{parsed\_data} (profil du candidat structuré)
        est stocké en \texttt{JSONB}, ce qui permet des requêtes sur les sous-champs.
  \item \textbf{MongoDB} : schéma trop flexible, difficile à maintenir avec des relations
        (candidat ↔ offre ↔ match).
  \item Port utilisé : \textbf{5433} (non-standard pour éviter les conflits avec
        d'autres instances locales).
\end{itemize}
\end{whybox}

\textbf{Tables principales :}

\begin{center}
\begin{tabularx}{\textwidth}{lX}
\toprule
\textbf{Table} & \textbf{Description} \\
\midrule
\texttt{users} & Comptes (candidat / recruteur / admin), auth JWT + OAuth \\
\texttt{candidates} & Profil structuré du candidat + \texttt{parsed\_data} JSONB \\
\texttt{cv\_documents} & Fichier CV uploadé (référence vers stockage) \\
\texttt{job\_offers} & Offres d'emploi (titre, description, compétences requises) \\
\texttt{matches} & Score de matching candidat↔offre (0–100\%) \\
\texttt{notifications} & Alertes email (nouvelles candidatures, statut) \\
\texttt{access\_logs} & Journal RGPD — chaque consultation de données personnelles \\
\bottomrule
\end{tabularx}
\end{center}

\begin{infobox}[Fichiers clés]
  \filepath{backend/app/models/candidate.py} — modèle Candidate (\texttt{parsed\_data} JSONB).\\
  \filepath{backend/app/models/match.py} — table des scores de matching.\\
  \filepath{backend/app/models/access\_log.py} — journal RGPD.\\
  \filepath{backend/alembic/versions/} — 9 fichiers de migration.
\end{infobox}

\section{Frontend — React 18 + Vite}

\begin{whybox}[Pourquoi React et pas Vue ou Angular ?]
\begin{itemize}
  \item \textbf{Composants réutilisables} : le composant \texttt{MatchReport} est
        utilisé à la fois dans le dashboard recruteur et la page candidat.
  \item \textbf{Hooks} (useState, useEffect, useContext) : gestion d'état simple
        sans Redux.
  \item \textbf{Vite} : serveur de développement ultra-rapide (HMR instantané).
  \item \textbf{Angular} : trop verbose, TypeScript obligatoire, courbe d'apprentissage.
\end{itemize}
\end{whybox}

\begin{infobox}[Point d'attention — Instance Axios centralisée]
  \filepath{frontend/src/services/api.js} — \textbf{toujours} importer depuis ce fichier.
  Il injecte automatiquement le JWT dans chaque requête et redirige vers \texttt{/login}
  en cas de 401. Contourner ce fichier casserait l'authentification.
\end{infobox}

\section{ORM — SQLAlchemy 2 + Alembic}

\begin{itemize}
  \item \textbf{SQLAlchemy 2} : ORM — on écrit du Python, pas du SQL brut.
        Prévient les injections SQL.
  \item \textbf{Alembic} : versioning du schéma de BDD. Chaque modification de
        modèle génère un fichier de migration dans
        \filepath{backend/alembic/versions/}.
\end{itemize}

% ============================================================
\chapter{Pipeline d'ingestion des CVs}
% ============================================================

Le pipeline complet se déclenche sur \texttt{POST /api/upload-cv}
(\filepath{backend/app/routes/cv.py}).

\section{Étape 1 — Extraction du texte}

Le module \filepath{backend/app/services/extraction/} dispatche selon le type de fichier :

\begin{center}
\begin{tabular}{lll}
\toprule
\textbf{Type} & \textbf{Outil} & \textbf{Fichier} \\
\midrule
PDF textuel & PyMuPDF / pypdf & \texttt{pdf\_extractor.py} \\
DOCX (Word) & python-docx & \texttt{word\_extractor.py} \\
Image / PDF scanné & EasyOCR & \texttt{ocr\_extractor.py} \\
\bottomrule
\end{tabular}
\end{center}

\begin{whybox}[Pourquoi EasyOCR et pas Tesseract ?]
\begin{itemize}
  \item EasyOCR utilise du \textbf{Deep Learning} (CRAFT + CRNN) — meilleure précision
        sur les documents bruités, inclinés ou avec des polices non standard.
  \item Tesseract est basé sur de l'OCR classique, plus fragile sur les CVs
        scannés de mauvaise qualité.
  \item EasyOCR supporte le français et l'arabe nativement.
\end{itemize}
\end{whybox}

\section{Étape 2 — Parsing NLP}

L'orchestrateur \filepath{backend/app/services/nlp/nlp\_parser.py} appelle
les extracteurs dans cet ordre :

\begin{enumerate}
  \item \textbf{ContactExtractor} — nom, email, téléphone, adresse.
  \item \textbf{EntityExtractor} — entités nommées via spaCy NER.
        Optionnel : \textbf{HFCamembertNameExtractor} si \texttt{HF\_ENABLE\_CAMEMBERT\_NAME=1}.
  \item \textbf{SkillsExtractor} — compétences techniques et soft skills.
        Fichier : \filepath{backend/app/services/nlp/skills\_extractor.py}
  \item \textbf{FormationExtractor} — diplômes, établissements, années.
        Fichier : \filepath{backend/app/services/nlp/formation\_extractor.py}
  \item \textbf{ExperienceExtractor} — postes, entreprises, durées.
        Fichier : \filepath{backend/app/services/nlp/experience\_extractor.py}
  \item \textbf{Langues} — langues détectées automatiquement (\texttt{\_extract\_langues}).
\end{enumerate}

\begin{whybox}[Pourquoi spaCy \texttt{fr\_core\_news\_md} et pas NLTK ?]
\begin{itemize}
  \item spaCy propose une \textbf{pipeline NLP complète} (tokenisation, POS, dépendances, NER)
        en un seul modèle.
  \item Le modèle \texttt{fr\_core\_news\_md} est entraîné sur du français journalistique
        et institutionnel — bien adapté aux CVs.
  \item NLTK est un toolkit académique sans modèle NER français de qualité.
  \item \textbf{Important} : le modèle n'est pas installable via pip. Il faut :
        \texttt{python -m spacy download fr\_core\_news\_md}
\end{itemize}
\end{whybox}

\section{Étape 3 — Stockage structuré}

Le résultat du parsing est stocké dans \texttt{Candidate.parsed\_data}
(champ JSONB PostgreSQL). Sa structure est définie dans
\filepath{backend/app/schemas/cv\_data.py}.

Exemple de structure :
\begin{lstlisting}[language=Python]
{
  "identite": {"nom_complet": "...", "email": "..."},
  "competences": {"techniques": ["React", "Python"], "soft": ["Leadership"]},
  "experiences": [{"poste": "Dev React", "duree_mois": 24, ...}],
  "formations": [{"diplome": "Licence Info", "etablissement": "Esprit"}],
  "secteur_detecte": {"secteur": "informatique", "label": "Développeur"},
  "metadata": {"spacy_fallback": false, "version": "2.1"}
}
\end{lstlisting}

% ============================================================
\chapter{Moteur de Matching IA}
% ============================================================

\section{Vue d'ensemble du scoring}

Le score final (0–100\%) est calculé par le scoreur principal
\filepath{backend/app/services/matching\_sandbox/bert\_scorer.py}
via une architecture à \textbf{3 niveaux} :

\begin{center}
\begin{tabular}{clc}
\toprule
\textbf{Niveau} & \textbf{Composant} & \textbf{Rôle} \\
\midrule
1 & BGE-M3 (bi-encoder) & Embeddings sémantiques \\
2 & Cross-Encoder (reranker) & Score sémantique précis \\
3 & MLP Fusion (7 features) & Score final combiné \\
\bottomrule
\end{tabular}
\end{center}

\section{Composant 1 — BGE-M3 (Bi-encoder)}

\begin{infobox}[Qu'est-ce que BGE-M3 ?]
BGE-M3 (\textbf{BAAI/bge-m3}) est un modèle de \textit{sentence embeddings}
multilingue développé par le Beijing Academy of Artificial Intelligence.
\begin{itemize}
  \item \textbf{568 millions} de paramètres, vecteurs de dimension \textbf{1024}.
  \item Supporte \textbf{100+ langues} — dont français, arabe et anglais en même temps.
  \item Basé sur l'architecture \textbf{XLM-RoBERTa} fine-tuné pour la similarité sémantique.
  \item Il encode séparément le CV et l'offre, puis calcule la \textbf{similarité cosinus}
        entre les deux vecteurs.
\end{itemize}
\end{infobox}

\begin{whybox}[Pourquoi BGE-M3 et pas CamemBERT ?]
\begin{itemize}
  \item \textbf{CamemBERT} est monolingue (français uniquement). Les CVs tunisiens
        contiennent souvent de l'anglais, du français et parfois de l'arabe.
  \item BGE-M3 est spécifiquement entraîné pour la \textbf{recherche sémantique} et
        la \textbf{similarité de documents}, pas juste pour le NER ou la classification.
  \item CamemBERT était utilisé dans les versions antérieures (Sprint 2-3).
        Le passage à BGE-M3 au Sprint 4 a amélioré les scores de matching.
\end{itemize}
\end{whybox}

\begin{codebox}[Localisation dans le code]
Fichier  : backend/app/services/matching_sandbox/bert_scorer.py
Ligne 734: _EMBEDDER_MODEL = "BAAI/bge-m3"
Ligne 758: class BERTMatchingScorer:
           -> méthode _encode() : calcule les embeddings
           -> méthode score() : retourne la similarité cosinus
\end{codebox}

\section{Composant 2 — Cross-Encoder (Reranker)}

\begin{infobox}[Qu'est-ce qu'un Cross-Encoder ?]
Contrairement au bi-encoder (BGE-M3 encode séparément les deux textes),
le Cross-Encoder (\textbf{BAAI/bge-reranker-v2-m3}) prend la \textbf{paire
(offre, CV) comme entrée unique}.
\begin{itemize}
  \item Il peut ainsi capturer des \textbf{interactions fines} entre les deux textes
        (un mot de l'offre qui répond à une phrase du CV).
  \item Résultat : un score scalaire unique, plus précis que la similarité cosinus.
  \item \textbf{Limitation} : beaucoup plus lent (pas de cache d'embeddings possible).
        C'est pourquoi il est utilisé en \textit{reranker} après un premier filtre BGE-M3.
\end{itemize}
\end{infobox}

\begin{codebox}[Localisation dans le code]
Fichier  : backend/app/services/matching_sandbox/bert_scorer.py
Ligne 737: _RERANKER_MODEL = "BAAI/bge-reranker-v2-m3"
Ligne 794: def _load_reranker(self):
           -> CrossEncoder chargé en lazy loading (au premier appel)
           -> max_length=512 pour correspondre au modèle
\end{codebox}

\section{Composant 3 — MLP Fusion (10 features)}

\subsection{Architecture du MLP}

Le MLP (\textit{Multi-Layer Perceptron}) est un réseau de neurones simple entraîné
\textbf{spécifiquement sur les données du projet}. La version actuelle (v3b1) prend
\textbf{10 features} en entrée, contre 7 dans la première version : trois signaux
métier supplémentaires ont été ajoutés (compétence du titre, criticité des
compétences, type d'offre) — détaillés au chapitre~\ref{chap:ameliorations}.

\begin{infobox}[Architecture — \filepath{backend/app/services/matching\_sandbox/bert\_scorer.py} (classe \texttt{\_ScoringMLPv3b1})]
$10\ \text{features} \rightarrow 64\ \text{neurones} \rightarrow 32\ \text{neurones} \rightarrow 1\ \text{score}$

Couches : Linéaire + BatchNorm + ReLU + Dropout 0.20, puis Linéaire + ReLU +
Dropout 0.10, puis Linéaire + Sigmoid (score entre 0 et 1).
\end{infobox}

\subsection{Les 10 features d'entrée}

\begin{center}
\begin{tabularx}{\textwidth}{clX}
\toprule
\textbf{\#} & \textbf{Feature} & \textbf{Description} \\
\midrule
1 & \texttt{sem\_sim} & Similarité sémantique BGE-M3 (bi-encoder) \\
2 & \texttt{skill\_rate\_w} & Ratio compétences pondéré par criticité (Phase 2) \\
3 & \texttt{exp\_score} & Score d'expérience (années vs. requis) \\
4 & \texttt{form\_score} & Score de formation, intégrant le domaine d'études (A3) \\
5 & \texttt{sem\_v2} & Score du Cross-Encoder (reranker) sur la paire \\
6 & \texttt{edu\_dom\_compat} & Compatibilité du domaine d'études candidat / offre \\
7 & \texttt{exp\_dom\_ratio} & Ratio d'expérience dans le domaine de l'offre \\
8 & \texttt{title\_skill\_match} & 1.0 si la compétence du titre est présente, 0.0 sinon (Phase 1) \\
9 & \texttt{crit\_skill\_ratio} & Ratio de compétences critiques matchées (Phase 2) \\
10 & \texttt{offer\_type\_enc} & Type d'offre encodé : 0=junior / 0.25=management / 0.5=tech / 1=senior (Phase 3) \\
\bottomrule
\end{tabularx}
\end{center}

\begin{whybox}[Pourquoi un MLP et pas une formule pondérée fixe ?]
\begin{itemize}
  \item Une formule fixe (ex: 40\% sémantique + 30\% skills + 30\% expérience) est
        \textbf{rigide} et ne s'adapte pas aux nuances métier.
  \item Le MLP \textbf{apprend les poids optimaux} depuis les données d'entraînement.
  \item Il peut apprendre des interactions non-linéaires (ex: un candidat sans
        diplôme requis mais avec 10 ans d'expérience est quand même pertinent).
\end{itemize}
\end{whybox}

\begin{infobox}[Méthode unique \texttt{\_compute\_mlp\_features()} — zéro décalage train/prod]
Le vecteur de 10 features est produit par \textbf{une seule méthode}, utilisée à la
fois pour l'inférence en production et pour la génération du dataset d'entraînement.
Cela garantit que le modèle apprend exactement sur les valeurs qu'il rencontrera en
production (pas de \textit{train/serving skew}).
\end{infobox}

\subsection{Entraînement du MLP}

\begin{infobox}[Fichiers d'entraînement (version v3b1)]
\begin{itemize}
  \item \filepath{backend/scripts/generate\_mlp\_dataset\_v3b1.py}
        → génère $\sim$360 paires (candidat, offre) couvrant 8 domaines ISCO et
        3 types d'offres, dont 30\% de cas hors-domaine. Encodage BGE-M3 en
        \textit{batch} (génération en moins de 6 secondes).
  \item \filepath{backend/scripts/train\_mlp\_v3b1.py}
        → entraîne le MLP $10\to64\to32\to1$ (split 80/10/10, early stopping).
        Sauvegarde les poids dans \filepath{data/models/fusion\_mlp/fusion\_mlp\_v3b1.pt}.
        Résultats : précision binaire 97\%, MAE 5,9\%.
  \item Le modèle entraîné est chargé au démarrage du serveur via le lifespan handler
        (\filepath{backend/app/main.py}). Influence du MLP dans le score final : 40\%.
\end{itemize}
\end{infobox}

\section{Domaine ISCO-08 — EscoService}

\begin{infobox}[Qu'est-ce que l'ISCO-08 ?]
L'ISCO-08 (\textit{International Standard Classification of Occupations 2008}) est
la classification internationale des métiers de l'OIT (Organisation Internationale
du Travail). Elle couvre \textbf{53 domaines professionnels}.
\end{infobox}

\textbf{Rôle dans TalentMatch :} avant de calculer le score, le système détecte
le domaine professionnel du CV et de l'offre via
\filepath{backend/app/services/matching\_sandbox/esco\_service.py}.

\begin{itemize}
  \item \textbf{Même domaine} (ex: Dev React ↔ Offre React) → pas de cap, score libre.
  \item \textbf{Domaines adjacents} (ex: Data Science ↔ Dev Backend) → cap à 0.65 (65\%).
  \item \textbf{Domaines différents} (ex: Médecin ↔ Offre Dev) → cap à 0.35 (35\%).
\end{itemize}

\begin{codebox}[Localisation]
Fichier : backend/app/services/matching_sandbox/esco_service.py
Données : data/esco/isco08_fr.json (53 groupes, keywords, skills_fr)
Méthode : EscoService.get_score_cap(domain1, domain2) -> float | None
\end{codebox}

\section{Score final — Pipeline complet}

\begin{center}
\begin{tabular}{cll}
\toprule
\textbf{Étape} & \textbf{Opération} & \textbf{Résultat} \\
\midrule
1 & BGE-M3 encode CV + Offre séparément & Vecteurs 1024-dim \\
2 & Similarité cosinus entre les deux vecteurs & Score sémantique brut \\
3 & Cross-Encoder (CV, Offre) en entrée unique & Score sémantique raffiné \\
4 & Calcul des 10 features (\texttt{\_compute\_mlp\_features}) & Vecteur 10 valeurs \\
5 & MLP Fusion (40\%) + formule calibrée (60\%) & Score 0–1 \\
6 & Multiplicateur de domaine ISCO-08 + cap titre (Phase 1) & Score ajusté \\
7 & Multiplication × 100 & Score \% affiché \\
\bottomrule
\end{tabular}
\end{center}

% ============================================================
\chapter{Améliorations avancées du moteur de matching}
\label{chap:ameliorations}
% ============================================================

Ce chapitre décrit les améliorations apportées au moteur après la première version.
Elles se répartissent en deux familles : la \textbf{qualité des données d'entrée}
(phases A1–A3) et la \textbf{logique métier de scoring} (phases 1–3), suivies de la
boucle d'apprentissage humain et de l'évaluation systématique du modèle.

\section{Qualité des données — Ontologie de synonymes (A1)}

\begin{whybox}[Le problème résolu]
La première version reconnaissait mal les synonymes de compétences. «~ReactJS~»,
«~React.js~» et «~React~» n'étaient pas reliés ; «~recruiting~» (anglais) ne
matchait pas «~recrutement~». Pire, le dictionnaire généré automatiquement
contenait du bruit (ex: «~react~»~$\rightarrow$~«~nuclear energy~»).
\end{whybox}

Une ontologie propre a été construite dans \filepath{data/skills/skills\_aliases.json}
(\textbf{5\,047 compétences} normalisées, FR + EN, multi-domaine). Elle fusionne :
un dictionnaire IT curé manuellement, un dictionnaire métier (Finance, RH, Santé,
Droit, BTP, Logistique\ldots), et les labels ESCO officiels. La normalisation
(minuscules, suppression des accents) est appliquée \textbf{avant} le calcul du
ratio de compétences.

\begin{codebox}[Génération]
Script : scripts/build_skills_aliases.py
Sortie : data/skills/skills_aliases.json (5047 entrées)
Test   : "ReactJS" == "React.js" == "React" -> match (8/8 validations)
\end{codebox}

\section{Qualité des données — Extraction multi-domaine (A2, A3)}

\textbf{A2 — Compétences non-IT.} L'extracteur de compétences
(\filepath{skills\_extractor.py}) a été corrigé pour gérer les compétences
composées («~contrôle de gestion~», «~gestion prévisionnelle des emplois~») et les
domaines non informatiques. Validation : \textbf{100\%} de rappel sur des CV de
référence en Comptabilité, RH et Santé (38/38 compétences extraites).

\textbf{A3 — Domaine d'études.} L'extracteur de formations
(\filepath{formation\_extractor.py}) ne se limite plus au niveau du diplôme
(Bac+5) mais en extrait aussi le \textbf{domaine} (informatique, finance,
droit\ldots). Ainsi «~Master Informatique~» et «~Master Finance~», de même niveau,
sont désormais distingués pour le scoring de la formation.

\section{Logique métier — Compétence centrale du titre (Phase 1)}

\begin{whybox}[Le problème résolu]
Un développeur Java avec 4 ans d'expérience obtenait un meilleur score qu'un
développeur Python junior pour un poste «~Développeur Python~» — alors que Python
est la compétence \textbf{obligatoire} indiquée dans le titre.
\end{whybox}

Le système détecte la compétence centrale du titre de manière \textbf{sémantique}
(via BGE-M3, sans liste codée en dur) : il compare le titre à chaque compétence
requise et retient la plus proche. Si le candidat ne possède pas cette compétence
\textbf{et} qu'il est dans le même domaine, son score est plafonné
($\times\,0{,}50$). Résultat : le développeur Java passe de 69\% à 34\% pour un
poste Python.

\section{Logique métier — Pondération par criticité (Phase 2)}

Toutes les compétences n'ont pas la même importance. Le système les classe en
trois niveaux de poids :

\begin{center}
\begin{tabular}{lcl}
\toprule
\textbf{Niveau} & \textbf{Poids} & \textbf{Exemple (poste Dev Python)} \\
\midrule
Critique (compétence du titre) & $\times 3$ & Python \\
Requise (métier)               & $\times 2$ & FastAPI, PostgreSQL \\
Générique (outil transversal)  & $\times 1$ & Git, Docker, Excel \\
\bottomrule
\end{tabular}
\end{center}

Le ratio de compétences est alors normalisé par la somme des poids, ce qui donne
plus de valeur aux compétences réellement discriminantes.

\section{Logique métier — Recalibration par type d'offre (Phase 3)}

Les poids des quatre dimensions (compétences, sémantique, expérience, formation)
ne sont plus fixes : ils s'adaptent au type d'offre détecté.

\begin{center}
\begin{tabular}{lcccc}
\toprule
\textbf{Type d'offre} & \textbf{Skills} & \textbf{Sém.} & \textbf{Exp.} & \textbf{Form.} \\
\midrule
Tech spécialisé   & 55\% & 20\% & 15\% & 10\% \\
Senior 5+ ans     & 40\% & 20\% & 30\% & 10\% \\
Stage / Junior    & 45\% & 25\% & 5\%  & 25\% \\
Management / RH   & 35\% & 25\% & 25\% & 15\% \\
\bottomrule
\end{tabular}
\end{center}

Le seuil de décision «~Entretien recommandé~» s'ajuste lui aussi (65\% pour un
poste junior, 70\% pour un senior).

\section{Apprentissage humain dans la boucle}

Pour ne pas apprendre uniquement sur une formule (raisonnement circulaire), un outil
d'annotation (\filepath{scripts/label\_pairs.py}) permet à un recruteur d'attribuer
un label réel (0 = mauvais, 0{,}5 = moyen, 1 = bon) à des paires (offre, CV)
sélectionnées de façon stratifiée. Ces labels sont stockés dans
\filepath{data/human\_labels.csv} et mélangés au dataset formule lors du
réentraînement, avec un \textbf{poids double} pour la vérité terrain humaine
(\filepath{train\_mlp\_human\_labels.py}). Un premier lot de 40 annotations a
confirmé un accord de 39/40 avec le modèle, validant sa calibration.

\section{Évaluation systématique du modèle}

Un \textit{benchmark} (\filepath{scripts/benchmark\_model.py}) mesure les
performances du modèle par \textbf{catégorie} de situation, pour cartographier
précisément ses forces et faiblesses.

\begin{center}
\begin{tabularx}{\textwidth}{Xcl}
\toprule
\textbf{Catégorie testée} & \textbf{Résultat} & \textbf{Statut} \\
\midrule
Match parfait (même domaine)        & 3/3 & Fort \\
Rejet hors-domaine total            & 3/3 & Fort \\
Domaines adjacents (zone grise)     & 2/2 & Fort \\
Synonymes EN/FR (\textit{aliases})  & 2/2 & Fort \\
Cap compétence titre (Phase 1)      & 2/2 & Fort \\
Surqualification                    & 1/1 & Fort \\
Sous-qualification expérience       & 2/2 & Fort \\
Discrimination junior (exp. faible) & 1/2 & Moyen \\
\midrule
\textbf{Score global}               & \textbf{16/17} & \textbf{94\%} \\
\bottomrule
\end{tabularx}
\end{center}

\textbf{Force principale.} Le modèle distingue très nettement les bons profils
(85–91\%) des profils hors-domaine (5–18\%), y compris sur les métiers non
informatiques, et reconnaît les synonymes multilingues.

\textbf{Point faible identifié.} Sur les offres junior/stage, le poids élevé de la
formation (25\%) peut gonfler le score d'un candidat ayant le bon diplôme mais peu de
compétences. Cet axe d'amélioration est documenté pour une itération future.

\begin{infobox}[Tests de non-régression]
En complément du benchmark, deux suites garantissent la stabilité :
\filepath{test\_matching\_reel.py} (15 cas de scoring, dont 6 multi-domaine) et
\filepath{test\_ranking.py} (classement validé par corrélation de Spearman sur
3 scénarios : $\rho = 0{,}90$ à $1{,}00$).
\end{infobox}

% ============================================================
\chapter{Sécurité et conformité RGPD}
% ============================================================

\section{Authentification JWT}

\begin{itemize}
  \item \textbf{JSON Web Token} (JWT) : standard RFC 7519. Un token signé est
        émis à la connexion et doit être envoyé dans chaque requête
        (header \texttt{Authorization: Bearer <token>}).
  \item Le token contient l'ID utilisateur et le rôle (candidat / recruteur / admin).
  \item Expiration configurable via \texttt{ACCESS\_TOKEN\_EXPIRE\_MINUTES} dans
        \filepath{backend/.env}.
  \item Implémentation : \filepath{backend/app/services/auth\_service.py}
\end{itemize}

\begin{whybox}[Pourquoi JWT et pas sessions classiques ?]
\begin{itemize}
  \item \textbf{Stateless} : le serveur n'a pas besoin de stocker les sessions en BDD.
        Scalable horizontalement (plusieurs instances du serveur).
  \item \textbf{Déployable sur Railway} sans partage de session entre instances.
\end{itemize}
\end{whybox}

\section{Rôles et contrôle d'accès}

Trois rôles définis dans \filepath{backend/app/dependencies.py} :

\begin{center}
\begin{tabular}{ll}
\toprule
\textbf{Rôle} & \textbf{Accès} \\
\midrule
Candidat & Upload CV, voir ses candidatures, liste des offres \\
Recruteur & Gérer ses offres, voir les résultats de matching \\
Admin & Tout + gestion utilisateurs + logs RGPD \\
\bottomrule
\end{tabular}
\end{center}

\section{Conformité RGPD}

\begin{itemize}
  \item \textbf{Consentement explicite} : lors de la candidature, le candidat coche
        une case RGPD. Visible sur la capture \textit{postuleCv.png}.
  \item \textbf{Journal d'accès} : table \texttt{access\_logs}
        (\filepath{backend/app/models/access\_log.py}) — chaque consultation de
        données personnelles est enregistrée (qui, quand, quoi).
  \item \textbf{Droit à l'effacement} : endpoint dédié de suppression complète.
  \item \textbf{Conservation limitée} : suppression automatique après 2 ans d'inactivité.
\end{itemize}

% ============================================================
\chapter{Interface utilisateur (Frontend)}
% ============================================================

\section{Architecture React}

\begin{infobox}[Fichier central — \filepath{frontend/src/services/api.js}]
Toutes les requêtes API passent par cette instance Axios unique.
Elle injecte automatiquement le JWT, gère le timeout (8s) et redirige
vers \texttt{/login} si le serveur répond 401.
\textbf{Ne jamais appeler Axios directement} dans un composant — toujours
passer par ce fichier.
\end{infobox}

\section{Pages principales}

\begin{center}
\begin{tabularx}{\textwidth}{lX}
\toprule
\textbf{Page / Fichier} & \textbf{Description} \\
\midrule
\filepath{frontend/src/pages/Dashboard.jsx} & Tableau de bord recruteur \\
\filepath{frontend/src/pages/OffersList.jsx} & Liste des offres (côté candidat) \\
\filepath{frontend/src/pages/OfferDetail.jsx} & Détail + bouton Postuler \\
\filepath{frontend/src/pages/OfferNew.jsx} & Création d'une nouvelle offre \\
\filepath{frontend/src/components/MatchReport.jsx} & Rapport de matching IA détaillé \\
\bottomrule
\end{tabularx}
\end{center}

\section{Rapport de matching — MatchReport}

Le composant \filepath{frontend/src/components/MatchReport.jsx} affiche pour chaque
candidat :
\begin{itemize}
  \item Score global (0–100\%) avec code couleur.
  \item Décomposition par dimension (Compétences, Expérience, Formation, Sémantique).
  \item Compétences présentes vs manquantes.
  \item Bouton « Analyse IA (Mistral) » → rapport narratif généré par LLM.
\end{itemize}

% ============================================================
\chapter{Méthodologie — Agile Scrum}
% ============================================================

\section{Organisation des sprints}

Le projet est organisé en \textbf{6 sprints d'1 mois} de Février à Août 2026.

\begin{center}
\begin{tabularx}{\textwidth}{cllX}
\toprule
\textbf{Sprint} & \textbf{Mois} & \textbf{Statut} & \textbf{Livrables principaux} \\
\midrule
1 & Février 2026 & \textcolor{codegreen}{\textbf{OK}} & CDC, conception UML, architecture, setup \\
2 & Mars 2026    & \textcolor{codegreen}{\textbf{OK}} & Upload CV, extraction texte, OCR, BDD \\
3 & Avril 2026   & \textcolor{codegreen}{\textbf{OK}} & Pipeline NLP spaCy, extraction skills \\
4 & Mai 2026     & \textcolor{orange}{\textbf{En cours}} & BGE-M3 + Cross-Encoder, matching, dashboard \\
5 & Juin 2026    & \textcolor{gray}{\textbf{TODO}} & Sécurité, tests intégration, notifications \\
6 & Juil–Août 2026 & \textcolor{gray}{\textbf{TODO}} & Déploiement Railway, documentation, soutenance \\
\bottomrule
\end{tabularx}
\end{center}

\section{Outils utilisés}

Git (GitHub), VS Code, PyCharm, Postman (tests API), Trello (backlog).
Reviews régulières avec Mme Sirine Nasri (encadrante société) et
Mme Loujayne Bouzarti (encadrante académique).

% ============================================================
\chapter{Avancement actuel — Sprint 4}
% ============================================================

\section{Fonctionnalités livrées}

\begin{itemize}
  \item[\checkmark] Authentification complète (JWT, OAuth Google/LinkedIn, OTP email).
  \item[\checkmark] Upload et extraction multi-format (PDF textuel, DOCX, PDF scanné via OCR).
  \item[\checkmark] Pipeline NLP complet : compétences, expériences, formations, langues.
  \item[\checkmark] Moteur de matching BGE-M3 + Cross-Encoder + MLP Fusion.
  \item[\checkmark] Dashboard recruteur : liste des offres, classement des candidats.
  \item[\checkmark] Interface candidat : liste des offres compatibles, candidature avec RGPD.
  \item[\checkmark] Score détaillé par dimension dans MatchReport.
  \item[\checkmark] Résumé IA du profil (Claude/Mistral).
  \item[\checkmark] Notifications (base implémentée).
\end{itemize}

\section{Démo — Scénario de test}

\begin{infobox}[Comptes de démo]
\begin{itemize}
  \item Recruteur : \texttt{selim@esprit.tn} / \texttt{Demo3S2026!}
  \item Admin : \texttt{admin1@esprit.tn} / \texttt{Demo3S2026!}
  \item Offre de démo : \textbf{Developpeur React Senior}
  \item 8 candidats avec scores de 93.7\% (Ahmed Benali) à 5.3\% (Amina Diallo)
\end{itemize}
\end{infobox}

% ============================================================
\chapter{Perspectives}
% ============================================================

\begin{enumerate}
  \item \textbf{Notifications email recruteur} — alerte à chaque nouvelle candidature.
  \item \textbf{Planification d'entretiens} — créneaux proposés par le recruteur, acceptés par le candidat.
  \item \textbf{Page candidat enrichie} — compétences manquantes identifiées, conseils personnalisés.
  \item \textbf{Amélioration continue par feedback} — le recruteur note « embauché / refusé »,
        ces données alimentent un réentraînement du MLP.
  \item \textbf{Génération de rapports via LLM (Mistral AI)} — rapport narratif complet sur chaque candidat.
  \item \textbf{Déploiement Railway} — conteneurisation Docker, CI/CD GitHub Actions.
\end{enumerate}

% ============================================================
\chapter*{Conclusion}
\addcontentsline{toc}{chapter}{Conclusion}
% ============================================================

TalentMatch apporte une réponse technique complète à la problématique du matching
CV / offres d'emploi. Le pipeline IA (BGE-M3 + Cross-Encoder + MLP Fusion) dépasse
les approches par mots-clés des solutions existantes en comprenant la \textbf{sémantique}
des profils et des offres.

Les choix technologiques (FastAPI, PostgreSQL, React, spaCy) ont été guidés par des
critères clairs : performance, maintenabilité et adéquation avec le périmètre du projet.
L'architecture modulaire (\textit{pipeline-shaped}) facilite les évolutions futures.

Le Sprint 4 marque le cœur du projet avec la livraison du moteur de matching IA.
Les sprints 5 et 6 compléteront la solution avec la sécurité renforcée, les tests
d'intégration et le déploiement en production sur Railway.

\vspace{1cm}
\begin{center}
{\color{navyblue}\rule{0.5\textwidth}{0.5pt}}\\[0.3cm]
\textit{Youssef Gara — Stage PFE 3S TalentMatch — Juin 2026}
\end{center}

\end{document}
